% GENERATED FILE. Edit canonical lessons, then run npm run generate:books.
% canonical-source-hash: fnv1a64:656932f09f53ef1e
% canonical-lessons: ES-C470-paraguas, ES-C470-mango, ES-C470-encontrar, ES-C470-avisar, ES-R470-repaso-paraguas, ES-C470-sintesis-paraguas

\chapter{Un paraguas azul}
\label{ch:es-paraguas}
\hlchaptermodality{470}

\begin{chapteropening}
\textbf{By the end of this chapter:} \emph{I can deal with something I have lost: read a message describing a mislaid object by colour and by the part you hold, say that I cannot find it anywhere, and ask to be told rather than ask anyone to go looking.}

Read a message to the building group about a mislaid umbrella, then answer it both ways.
\end{chapteropening}

\section[el paraguas]{el paraguas — the umbrella}

\label{lesson:ES-C470-paraguas}



\begin{quote}
Say \textbf{la parada} and \textbf{el cumpleaños}. One lends this word its first half, the other lends it its shape.
\end{quote}

\subsection*{You'll want to know: el paraguas}
\textbf{el paraguas} — the umbrella.

\begin{quote}
Me he dejado el \textbf{paraguas} en el ascensor.
\end{quote}

Look at it before you are told anything: \textbf{para} plus \textbf{aguas}. A verb and a plural noun fused into one word, which is the move \emph{el cumpleaños} taught you — \emph{cumple} plus \emph{años}, a sentence that hardened.

And like \emph{cumpleaños}, it wears a plural \textbf{-s} and is singular: \textbf{el} paraguas, one umbrella.

\begin{grammarlens}[title={Grammar Lens: the plural that does not move}]
So what is the plural? Nothing changes.

\noindent
\begin{tabularx}{\linewidth}{@{}>{\raggedright\arraybackslash}X>{\raggedright\arraybackslash}X@{}}
\toprule
\textbf{one} & \textbf{more than one} \\
\midrule
\textbf{el} paraguas & \textbf{los} paraguas \\
\bottomrule
\end{tabularx}

This is not a quirk of this word. El análisis gave you the rule: a word already ending in \textbf{s}, with its stress somewhere other than the last syllable, takes no plural ending. \emph{El lunes} / \emph{los lunes} works the same way.

The article is doing all the counting. Drop it and the sentence cannot tell you how many umbrellas there are.
\end{grammarlens}

\begin{cousinweb}
The \textbf{para-} is the one \emph{la parada} already unpacked: Latin \textbf{parare}, to \textbf{make ready} — and \emph{parada} reminded you it is not the Greek \emph{para-} of \emph{parallel}. A \textbf{parasol} is made ready against the sun, a \textbf{parachute} against a fall. A \textbf{paraguas} is made ready against the waters.

Which is worth a second look, because English named the same object from the opposite direction:

\noindent
\begin{tabularx}{\linewidth}{@{}>{\raggedright\arraybackslash}X>{\raggedright\arraybackslash}X@{}}
\toprule
\textbf{language} & \textbf{what the word notices} \\
\midrule
Spanish \textbf{paraguas} & the \textbf{water} it keeps off \\
English \textbf{umbrella} & the \textbf{shade} it makes \\
\bottomrule
\end{tabularx}

\emph{Umbrella} is Latin \emph{umbra}, a shadow, with a diminutive on it — a little shade, and English took it from Italian. So one language named the object for the shade it casts and the other for the water it turns away, and both are describing the same thing held over the same head.
\end{cousinweb}

\subsection*{Your turn}
Say these aloud:

\begin{itemize}
\raggedright
  \item I have left the umbrella in the lift
  \item the plural of \emph{el paraguas}
  \item what \emph{parare} means
\end{itemize}

\begin{tcolorbox}[breakable,colback=teal!4,colframe=teal!35!black,title={Before you move on}]
The umbrella? (\textbf{El paraguas}.) And two of them? (\textbf{Los paraguas} — the word does not change.) What is it made ready against? (\textbf{The waters}.)
\end{tcolorbox}
\section[el mango]{el mango — the handle}

\label{lesson:ES-C470-mango}



\begin{quote}
Say \textbf{la mano} and \textbf{el paraguas}. This word is a relative of the first and a part of the second.
\end{quote}

\subsection*{You'll want to know: el mango}
\textbf{el mango} — the handle.

\begin{quote}
Un paraguas azul, con el \textbf{mango} de madera.
\end{quote}

The handle of a knife, a pan, a door, an umbrella: whatever part of a thing is built for the hand.

Which is what the word says. It is a cousin of \textbf{la mano}, and once you hear that, you will not need to store it separately.

\begin{grammarlens}[title={Grammar Lens: two mangos, and only one of them grows}]
There is a second \textbf{mango} in Spanish, and it is a completely different word that happens to be spelled the same. This is the situation la cola prepared you for.

\noindent
\begin{tabularx}{\linewidth}{@{}>{\raggedright\arraybackslash}X>{\raggedright\arraybackslash}X@{}}
\toprule
\textbf{the word} & \textbf{where it comes from} \\
\midrule
\textbf{el mango} — handle & Latin \emph{manus}, the hand \\
\textbf{el mango} — the fruit & Portuguese \emph{manga}, from Malay \\
\bottomrule
\end{tabularx}

The Academy traces the fruit through Portuguese \textbf{manga} to Malay \emph{mangga}; behind that stands a Dravidian word, Tamil \emph{māṅkāy}. It travelled by ship and the handle did not. Nothing connects them but the spelling.

Context settles it every time, so this is not a hazard — it is a reminder that \emph{same letters} and \emph{same word} are different claims.
\end{grammarlens}

\begin{cousinweb}
Latin \textbf{manus}, the hand — the root la mano already gave you, along with \emph{manual}, \emph{manufacture} and \emph{maintain}.

A \textbf{mango} is the hand-part. Spanish built it on a late Latin form of that root and let the meaning settle on the graspable end of a thing.

English does not have this word, so it has nothing to lend you here. What it does have is \emph{manual}, \emph{manuscript}, \emph{manacle} — all of them hands. Put \emph{mango} beside them and it stops looking like a fruit.
\end{cousinweb}

\subsection*{Your turn}
Say these aloud:

\begin{itemize}
\raggedright
  \item a blue umbrella with a wooden handle
  \item the two mangos, and where each comes from
  \item which Latin word is behind the handle
\end{itemize}

\begin{tcolorbox}[breakable,colback=teal!4,colframe=teal!35!black,title={Before you move on}]
The handle? (\textbf{El mango}.) Which body part is inside it? (\textbf{La mano} — Latin \emph{manus}.) And the fruit? (\textbf{A different word entirely.})
\end{tcolorbox}
\section[encontrar]{encontrar — to find}

\label{lesson:ES-C470-encontrar}



\begin{quote}
Say \textbf{buscar}. You have had the looking for a long time. Say \textbf{el mango}. You still have no word for the moment the looking succeeds.
\end{quote}

\subsection*{You'll want to know: encontrar}
\textbf{encontrar} — to find.

\begin{quote}
No lo \textbf{encuentro} por ninguna parte.
\end{quote}

\emph{Buscar} and \emph{encontrar} are the two halves of one search, and Spanish keeps them strictly apart:

\noindent
\begin{tabularx}{\linewidth}{@{}>{\raggedright\arraybackslash}X>{\raggedright\arraybackslash}X@{}}
\toprule
\textbf{what you do} & \textbf{what happens} \\
\midrule
\textbf{buscar} — look for & the effort \\
\textbf{encontrar} — find & the result \\
\bottomrule
\end{tabularx}

\emph{Busco mi paraguas} is work in progress. \emph{He encontrado mi paraguas} is over. You can \emph{buscar} all afternoon and never \emph{encontrar}.

\begin{grammarlens}[title={Grammar Lens: the o that breaks}]
You wrote \emph{encuentro} above, not \emph{encontro}. That is mostrar's rule arriving again, unchanged: push the stress back onto the stem and the \textbf{o} breaks into \textbf{ue}.

\noindent
\begin{tabularx}{\linewidth}{@{}>{\raggedright\arraybackslash}X>{\raggedright\arraybackslash}X@{}}
\toprule
\textbf{stress on the stem} & \textbf{stress on the ending} \\
\midrule
enc\textbf{ue}ntro, enc\textbf{ue}ntras & encontr\textbf{a}mos \\
enc\textbf{ue}ntra, enc\textbf{ue}ntran & encontr\textbf{á}is \\
\bottomrule
\end{tabularx}

The infinitive keeps its \textbf{o} because the stress is on the ending there — \emph{encon-TRAR}. Nothing about the verb is irregular once you know where the weight falls.
\end{grammarlens}

\begin{cousinweb}
\textbf{en-} plus \textbf{contra}: to come up \textbf{against} something.

That is the whole picture, and it is more honest than the English word. \emph{Find} says nothing about how. \emph{Encontrar} says you went along and the thing was in your way.

English has the same word by a different road: \textbf{encounter}, borrowed from French, is \emph{in} plus \emph{contra} exactly. To encounter somebody is to come against them — which is why an encounter can be a meeting or a fight, and \emph{un encuentro} in Spanish is both as well.

So \emph{encontrarse con alguien} is to run into somebody. The \emph{contra} has not gone anywhere.
\end{cousinweb}

\subsection*{Your turn}
Say these aloud:

\begin{itemize}
\raggedright
  \item I can't find it anywhere
  \item I am looking for it, and I have not found it
  \item what \emph{contra} is doing inside the word
\end{itemize}

\begin{tcolorbox}[breakable,colback=teal!4,colframe=teal!35!black,title={Before you move on}]
To find? (\textbf{Encontrar}.) I find? (\textbf{Encuentro} — the stress moved onto the stem.) What is \emph{contra} doing there? (\textbf{The thing came up against you}.)
\end{tcolorbox}
\section[avisar]{avisar — to let know}

\label{lesson:ES-C470-avisar}



\begin{quote}
Say \textbf{el aviso}. Say \textbf{encontrar}. If somebody finds the thing, this is what you want them to do next.
\end{quote}

\subsection*{You'll want to know: avisar}
\textbf{avisar} — to let know, to notify.

\begin{quote}
Si lo habéis visto, \textbf{avisadme}, por favor.
\end{quote}

You have met this verb once already without being given it. El aviso named it in passing — \emph{\textquotedblleft{}behind it is the verb avisar, to notify\textquotedblright{}} — and then went back to the noun. Here it is as yours.

\emph{Avísame} — let me know. \emph{Avisar a alguien} — to tell somebody. \emph{Avisar de algo} — to give notice of something.

\begin{grammarlens}[title={Grammar Lens: it shows, it does not order}]
\emph{El aviso} made a point that the verb inherits whole: an \emph{aviso} is not an instruction but a \textbf{making-visible}. Somebody puts a thing where you can see it, and what you do next is yours.

So \emph{avisar} is the weakest of the telling verbs, and that is its use.

\noindent
\begin{tabularx}{\linewidth}{@{}>{\raggedright\arraybackslash}X>{\raggedright\arraybackslash}X@{}}
\toprule
\textbf{you want} & \textbf{you say} \\
\midrule
somebody to know, no more & \textbf{avisar} \\
somebody to do it & \emph{decir} that they must \\
\bottomrule
\end{tabularx}

\emph{Avísame si lo encuentras} asks for information and nothing else. It is what you say about a paraguas you left behind: you are not giving an order to a neighbour, you are asking to be shown something.
\end{grammarlens}

\begin{cousinweb}
\emph{El aviso} already traced this: Latin \textbf{ad vīsum}, toward what has been seen, on \textbf{vidēre}, to see. There is nothing to add to the derivation, so take instead the thing it predicts.

Every relative of \emph{avisar} is about \textbf{seeing}, not speaking:

\begin{itemize}
\raggedright
  \item English \textbf{advise} is the same formation — to put something in somebody's view;
  \item \textbf{advice} is what you were shown;
  \item \textbf{evident}, \textbf{vision}, \textbf{provide}, \textbf{survey} are all seeings.
\end{itemize}

A word for telling, built entirely out of sight. That is why \emph{avisar} feels lighter than it looks: nobody in it is talking.
\end{cousinweb}

\subsection*{Your turn}
Say these aloud:

\begin{itemize}
\raggedright
  \item let me know if you find it
  \item I left it in the lift; tell me if you have seen it
  \item which Latin verb is underneath
\end{itemize}

\begin{tcolorbox}[breakable,colback=teal!4,colframe=teal!35!black,title={Before you move on}]
To let know? (\textbf{Avisar}.) Let me know? (\textbf{Avísame}.) Is it an order? (\textbf{No — it makes something visible}.)
\end{tcolorbox}
\section[(el paraguas, el mango, encontrar …]{(el paraguas, el mango, encontrar, avisar) — from cold}

\label{lesson:ES-R470-repaso-paraguas}



\begin{quote}
Four words, no prompts. The thing you left behind, the part of it you would describe, what you cannot do, and what you want the finder to do.
\end{quote}

\subsection*{You'll want to know: the four, cold}
\noindent
\begin{tabularx}{\linewidth}{@{}>{\raggedright\arraybackslash}X>{\raggedright\arraybackslash}X@{}}
\toprule
\textbf{English} & \textbf{Spanish} \\
\midrule
the umbrella & \textbf{el paraguas} \\
the handle & \textbf{el mango} \\
to find & \textbf{encontrar} \\
to let know & \textbf{avisar} \\
\bottomrule
\end{tabularx}

\begin{grammarlens}[title={Grammar Lens: one English verb, two Spanish ones}]
English says \emph{look for it and tell me}. Spanish makes you choose twice, and the choices are not the ones you would guess.

\noindent
\begin{tabularx}{\linewidth}{@{}>{\raggedright\arraybackslash}X>{\raggedright\arraybackslash}X@{}}
\toprule
\textbf{the English} & \textbf{what Spanish splits it into} \\
\midrule
look / find & \emph{buscar} is the effort, \textbf{encontrar} the result \\
tell & \textbf{avisar} shows; \emph{decir} says \\
\bottomrule
\end{tabularx}

\emph{Avísame si lo encuentras} holds both splits in five words: let me know — not order me, not explain to me — if the search actually lands.

That sentence is the whole chapter, and you could not have built it a lesson ago.
\end{grammarlens}

\begin{cousinweb}
Three of the four are older than they look, and each keeps a picture inside.

\noindent
\begin{tabularx}{\linewidth}{@{}>{\raggedright\arraybackslash}X>{\raggedright\arraybackslash}X@{}}
\toprule
\textbf{word} & \textbf{the picture} \\
\midrule
el paraguas & made \textbf{ready} against the waters \\
el mango & the part built for the \textbf{hand} \\
encontrar & to come up \textbf{against} a thing \\
\bottomrule
\end{tabularx}

\emph{Avisar} is the fourth, and its picture is \textbf{seeing} rather than telling — which is what made it the gentle verb in the lens above.
\end{cousinweb}

\subsection*{Your turn}
Say these aloud:

\begin{itemize}
\raggedright
  \item a blue umbrella with a wooden handle
  \item I can't find it anywhere, let me know if you see it
  \item which of the four is about seeing
\end{itemize}

\begin{tcolorbox}[breakable,colback=teal!4,colframe=teal!35!black,title={Before you move on}]
The umbrella? (\textbf{El paraguas}.) The handle? (\textbf{El mango}.) To find? (\textbf{Encontrar}.) To let know? (\textbf{Avisar}.)
\end{tcolorbox}
\section[Practice]{(el mensaje del vecino) — read it, then answer}

\label{lesson:ES-C470-sintesis-paraguas}



\begin{quote}
Say \textbf{el paraguas} and \textbf{avisar}. One is lost and the other is what the message is asking for.
\end{quote}

\begin{tcolorbox}[breakable,skin=enhanced,colback=orange!6,colframe=orange!45!black,boxrule=0.5pt,arc=1mm,left=6pt,right=6pt,top=4pt,bottom=4pt,fonttitle=\bfseries,title={Read this}]
To the building group, this morning:

\begin{quote}
Hola, ¿sois vosotros los del piso doce? Me he dejado un \textbf{paraguas} azul, con el \textbf{mango} de madera, en el ascensor esta mañana y ya no lo \textbf{encuentro} por ninguna parte. Si lo habéis visto, \textbf{avisadme}, por favor.
\end{quote}

Two things to notice before you answer.

\textbf{Azul, de madera.} The colour is the umbrella's and the wood is the handle's. Spanish hangs each description on the noun it belongs to, and \emph{de madera} is how you say what a thing is made of.

\textbf{Por ninguna parte.} The stacking is no … nada's rule, unchanged: \emph{no} in front of the verb and the negative word after it, both required. What is new is only the phrase. \emph{Ninguna parte} joins \emph{nada}, \emph{nadie} and \emph{nunca} as a negative word, and it is two words rather than one.

The trap is the English. \emph{No lo encuentro por ninguna parte} comes out as \emph{I can't find it \textbf{anywhere}}, which looks positive — so the pull is to reach for \emph{en cualquier parte}. Spanish keeps the negative on both halves.
\end{tcolorbox}

\begin{grammarlens}[title={Grammar Lens: what the message is actually asking}]
Read the last sentence again. It does not ask anyone to look.

\noindent
\begin{tabularx}{\linewidth}{@{}>{\raggedright\arraybackslash}X>{\raggedright\arraybackslash}X@{}}
\toprule
\textbf{what it says} & \textbf{what it does not say} \\
\midrule
\emph{si lo habéis visto} & \emph{buscadlo} \\
\bottomrule
\end{tabularx}

The writer is asking to be told about something that may already have happened, not asking the building to start a search. That is \emph{avisar} doing exactly the work it was built for, and it is the polite shape of this kind of message: you ask for information, not labour.

Answer in the same register.
\end{grammarlens}

\subsection*{Your turn}
Say what was lost, where, and what it looks like — colour and handle both.

Now answer the message twice. First: you have seen it, and you say where. Then: you have not, and you offer to let them know if it turns up.

\begin{tcolorbox}[breakable,colback=teal!4,colframe=teal!35!black,title={Before you move on}]
What is lost? (\textbf{Un paraguas azul con el mango de madera}.) What can the writer not do? (\textbf{Encontrarlo}.) What are they asking for? (\textbf{Que les avisen} — to be told.)
\end{tcolorbox}
